% Part: normal-modal-logic
% Chapter: filtrations
% Section: examples-of-filtrations

\documentclass[../../../include/open-logic-section]{subfiles}

\begin{document}

\olfileid{nml}{fil}{exf}

\olsection{Tuladha Filtrasi}

Kita durung nuduhake yen filtrasi pancen ana. Nanging kanggo saben
model $\mModel{M}$, ana akeh filtrasi saka $\mModel{M}$ liwat
$\Gamma$. Kita mbedakake mligine rong jinis: filtrasi paling alus lan
paling kasar. Filtrasi saka model kang padha bisa beda ing relasi
aksesibilitase (amarga \olref[fil]{defn:filtration} wis nemtokake
$W^*$ lan~$V^*$ kanthi langsung). Filtrasi paling alus nduweni
sesambungan antarjagad sithik-dhikite, dene filtrasi paling kasar
nduweni sesambungan akeh-akehe.

\begin{defn}
  Yen $\Gamma$ katutup tumrap subformula, filtrasi \emph{paling alus}
  $\mModel{M^*}$ saka model $\mModel{M}$ ditepakake kanthi:
  \[
  R^*[u][v] \quad \text{yen lan mung yen} \quad \exists u'\in [u] \;
  \exists v' \in [v] : Ru'v'.
  \]
\end{defn}

\begin{prop}\ollabel{prop:finest}
  Filtrasi paling alus $\mModel{M^*}$ pancen sawijining filtrasi.
\end{prop}

\begin{proof}
  Kita kudu mriksa yen $R^*$ kang ditepakake mangkono nyukupi
  \olref[fil]{defn:filtration}\olref[fil]{defn:filtration-R}.
  Kita priksa telung sarat iku siji-siji.

  Yen $Ruv$, amarga $u \in [u]$ lan $v \in [v]$, uga $R^*[u][v]$,
  mula \olref[fil]{defn:filtration-R1} kawujud.
  
  \iftag{prvBox}{\iftag{probBox}{Pamariksan
      \olref[fil]{defn:filtration-R2} kita tinggalake minangka latihan.}{Kanggo
        \olref[fil]{defn:filtration-R2}, upamane $\Box!A \in \Gamma$,
        $R^*[u][v]$, lan $\mSat{M}{\Box!A}[u]$. Miturut definisi
        $R^*$, ana $u' \equiv u$ lan $v' \equiv v$ kang
        $Ru'v'$. Amarga $u$ lan $u'$ sarujuk tumrap $\Gamma$, uga
        $\mSat{M}{\Box!A}[u']$, mula $\mSat{M}{!A}[v']$. Amarga
        $\Gamma$ katutup tumrap sub-!!{formula}s, $v$ lan $v'$ sarujuk
        tumrap $!A$, mula $\mSat{M}{!A}[v]$, kaya kang dikarepake.}}{}
  
  \iftag{prvDiamond}{\iftag{probDiamond}{Pamariksan
      \olref[fil]{defn:filtration-R3} kita tinggalake minangka latihan.}{Kanggo mriksa
        \olref[fil]{defn:filtration-R3}, upamane $\Diamond!A \in
        \Gamma$, $R^*[u][v]$, lan $\mSat{M}{!A}[v]$. Miturut definisi
        $R^*$, ana $u' \equiv u$ lan $v' \equiv v$ kang
        $Ru'v'$. Amarga $v$ lan $v'$ sarujuk tumrap~$\Gamma$, lan
        $\Gamma$ katutup tumrap sub-!!{formula}s, uga
        $\mSat{M}{!A}[v']$, mula $\mSat{M}{\Diamond !A}[u']$.
        Amarga $u$ lan $u'$ uga sarujuk tumrap $\Gamma$,
        $\mSat{M}{\Diamond !A}[u]$.}}{}
\end{proof}

\begin{probtag}{probBox,probDiamond}
  Jangkepna bukti \olref[nml][fil][exf]{prop:finest}.
\end{probtag}

\begin{defn}
  Yen $\Gamma$ katutup tumrap subformula, filtrasi \emph{paling kasar}
  $\mModel{M^*}$ saka model~$\mModel{M}$ ditepakake kanthi
  $R^*[u][v]$ yen lan mung yen
  \iftag{notprvBox,notprvDiamond}{sarat ing ngisor iki
    kawujud:}{\emph{loro-lorone} sarat ing ngisor iki kawujud:}
  \begin{tagenumerate}{prvBox,prvDiamond}
  \tagitem{prvBox}{\ollabel{defn:coarsest-Box}Yen $\Box!A \in \Gamma$
    lan $\mSat{M}{\Box!A}[u]$, mula
    $\mSat{M}{!A}[v]$\iftag{prvDiamond}{;}{.}}{}
  \tagitem{prvDiamond}{\ollabel{defn:coarsest-Diamond}Yen $\Diamond!A
    \in \Gamma$ lan $\mSat{M}{!A}[v]$, mula
    $\mSat{M}{\Diamond!A}[u]$.}{}
  \end{tagenumerate}
\end{defn}

\begin{prop}
  Filtrasi paling kasar $\mModel{M^*}$ pancen sawijining filtrasi.
\end{prop}

\begin{proof}
  Miturut definisi $R^*$, siji-sijine sarat kang isih kudu dipriksa
  yaiku implikasi saka $Ruv$ menyang $R^*[u][v]$. Dadi upamane
  $Ruv$. \iftag{prvBox}{Upamane $\Box!A \in \Gamma$ lan
    $\mSat{M}{\Box!A}[u]$; banjur cetha
    $\mSat{M}{!A}[v]$\iftag{prvDiamond}{, lan
      \olref{defn:coarsest-Box} kawujud. }{.}}{}\iftag{prvDiamond}{Upamane
    $\Diamond!A \in \Gamma$ lan $\mSat{M}{!A}[v]$. Banjur
    $\mSat{M}{\Diamond!A}[u]$ amarga~$Ruv$\iftag{prvBox}{, lan
      \olref{defn:coarsest-Diamond} kawujud}{}.}{}
\end{proof}

\begin{ex}
  Jupuken $W = \PosInt$, $Rnm$ yen lan mung yen $m = n + 1$, lan
  $V(p) = \Setabs{2n}{n \in \PosInt}$. Model
  $\mModel{M} = \tuple{W, R, V}$ digambarake ing
  \olref{fig:ex-filtration}. Jagad-jagade yaiku $1$, $2$, lan
  sateruse; saben jagad bisa ngakses persis siji jagad liyane---
  paneruse---lan $p$ bener ing kabeh lan mung wilangan genep.

  \begin{figure}
    \centering
    \begin{tikzpicture}[modal]
      \node[world] (1) [label=below:\mFalse{p}] {$1$};
      \node[world] (2) [label=below:\mTrue{p}, right=of 1] {$2$};
      \node[world] (3) [label=below:\mFalse{p}, right=of 2] {$3$};
      \node[world] (4) [label=below:\mTrue{p}, right=of 3] {$4$};
      \node[phantom] (5) [right=of 4] {};
      \draw[->] (1) to (2);
      \draw[->] (2) to (3);
      \draw[->] (3) to (4);
      \draw[dotted] (4) to (5);
    \end{tikzpicture}

    \begin{tikzpicture}[modal]
      \node[world] (1) [label=below:\mFalse{p}] {$[1]$};
      \node[world] (2) [label=below:\mTrue{p}, right=of 1] {$[2]$};
      \draw[->,bend left] (1) to (2);
      \draw[->,bend left] (2) to (1);
      
      \node[world] (3) [label=below:\mFalse{p},right=of 2] {$[1]$};
      \node[world] (4) [label=below:\mTrue{p}, right=of 3] {$[2]$};
      \draw[->,bend left] (3) to (4);
      \draw[->,bend left] (4) to (3);
      \draw[->,reflexive above] (4) to (4);
    \end{tikzpicture}
    \caption{Model tanpa wates lan filtrasi-filtrasine.}
    \ollabel{fig:ex-filtration}
  \end{figure}
  
  Saiki jupuken $\Gamma$ minangka himpunan sub-!!{formula}s saka
  $\Box p \lif p$, yaiku $\{p, \Box p, \Box p \lif p\}$. $p$ bener
  ing kabeh lan mung wilangan genep, $\Box p$ bener ing kabeh lan mung
  wilangan ganjil, mula $\Box p \lif p$ bener ing kabeh lan mung
  wilangan genep. Kanthi tembung liya, saben wilangan ganjil ndadekake
  $\Box p$ bener nanging $p$ lan $\Box p \lif p$ salah; saben wilangan
  genep ndadekake $p$ lan $\Box p \lif p$ bener, nanging $\Box p$
  salah. Dadi $W^* = \{ [1], [2] \}$, ing ngendi $[1] =
  \{1, 3, 5, \dots\}$ lan $[2] = \{2, 4, 6, \dots\}$. Amarga $2 \in
  V(p)$, $[2] \in V^*(p)$; amarga $1 \notin V(p)$, $[1] \notin
  V^*(p)$. Dadi $V^*(p) = \{[2]\}$.

  Saben filtrasi adhedhasar $W^*$ kudu nduweni relasi aksesibilitas
  kang ngemot $\tuple{[1], [2]}, \tuple{[2],[1]}$: amarga $R12$,
  kudu ana $R^*[1][2]$ miturut
  \olref[fil]{defn:filtration}\olref[fil]{defn:filtration-R1}, lan
  amarga $R23$ kudu ana $R^*[2][3]$, dene $[3]=[1]$. Relasi iku ora
  bisa ngemot $\tuple{[1],[1]}$: yen ngemot, kita bakal duwe
  $R^*[1][1]$, $\mSat{M}{\Box p}[1]$ nanging
  $\mSat/{M}{p}[1]$, kang nalisir karo
  \olref[fil]{defn:filtration-R2}. Ora ana sarat kang meksa utawa
  nglarang $R^*[2][2]$. Mula, ana rong filtrasi kang mungkin saka
  $\mModel{M}$, miturut rong relasi aksesibilitas iki:
  \[
  \{\tuple{[1],[2]}, \tuple{[2],[1]}\} \text{ lan }
  \{\tuple{[1],[2]}, \tuple{[2],[1]}, \tuple{[2],[2]}\}.
  \]
  Ing loro-lorone kasus, $p$ lan $\Box p \lif p$ salah dene $\Box p$
  bener ing~$[1]$; $p$ lan $\Box p \lif p$ bener dene $\Box p$ salah
  ing~$[2]$.
\end{ex}


\begin{prob}
  Gatekna model $\mModel{M} = \tuple{W, R, V}$ ing ngendi $W
  = \Setabs{0\sigma}{\sigma \in \Bin^*}$ iku himpunan rerangkening
  $0$s lan~$1$s kang diwiwiti~$0$, kanthi $R\sigma\sigma'$ yen lan
  mung yen $\sigma' = \sigma 0$ utawa $\sigma' = \sigma 1$, lan
  $V(p) = \Setabs{\sigma 0}{\sigma \in \Bin^*} \cap W$ lan
  $V(q) = \Setabs{\sigma 1}{\sigma \in \Bin^* \setminus \{1\}} \cap W$.
  Iki gambare:
  \begin{center}
    \begin{tikzpicture}[modal,
        node distance=2cm
      ]

      \node[world] (0) [label={below:\mTrue{p}\\\mFalse{q}},font=\tiny] {$0$};
      \node[world] (00) [label={below:\mTrue{p}\\\mFalse{q}},
        above right=of 0,font=\tiny] {$00$};
      \node[world] (000) [label={below:\mTrue{p}\\\mFalse{q}},
        right=of 00, yshift=.8cm, font=\tiny] {$000$};
      \node[phantom] (0000) [right=of 000, yshift=.5cm] {};
      \node[phantom] (0001) [right=of 000, yshift=-.5cm] {};
      \node[world] (001) [label={below:\mFalse{p}\\\mTrue{q}},
        right=of 00, yshift=-.8cm, font=\tiny] {$001$};
      \node[phantom](0010) [right=of 001, yshift=.5cm] {};
      \node[phantom](0011) [right=of 001, yshift=-.5cm] {};
      \node[world] (01) [label={below:\mFalse{p}\\\mTrue{q}},
        below right=of 0,font=\tiny] {$01$};
      \node[world] (010) [label={below:\mTrue{p}\\\mFalse{q}},
        right=of 01, yshift=.8cm,font=\tiny] {$010$};
      \node[phantom] (0100) [right=of 010, yshift=.5cm] {};
      \node[phantom] (0101) [right=of 010, yshift=-.5cm] {};
      \node[world] (011) [label={below:\mFalse{p}\\\mTrue{q}},
        right=of 01, yshift=-.8cm,font=\tiny] {$011$};
      \node[phantom] (0110) [right=of 011, yshift=.5cm] {};
      \node[phantom] (0111) [right=of 011, yshift=-.5cm] {};

      \draw[->] (0) to (00);
      \draw[->] (0) to (01);
      \draw[->] (00) to (000);
      \draw[dotted] (000) to (0000);
      \draw[dotted] (000) to (0001);
      \draw[->] (00) to (001);
      \draw[dotted] (001) to (0010);
      \draw[dotted] (001) to (0011);
      \draw[->] (01) to (010);
      \draw[dotted] (010) to (0100);
      \draw[dotted] (010) to (0101);
      \draw[->] (01) to (011);
      \draw[dotted] (011) to (0110);
      \draw[dotted] (011) to (0111);
    \end{tikzpicture}
  \end{center}
  Kita duwe $\mSat/{M}{\Box(p \lor q) \lif (\Box p \lor \Box q)}[w]$
  kanggo saben~$w$.

  Upamane $\Gamma$ himpunan sub-!!{formula}s saka~$\Box(p \lor q) \lif
  (\Box p \lor \Box q)$. Apa $W^*$ lan~$V^*$? Apa relasi
  aksesibilitas filtrasi paling alus saka~$\mModel{M}$? Kepiye
  filtrasi paling kasar?
\end{prob}

\end{document}
